\documentclass{article}
\usepackage{iclr2026_conference}
\input{math_commands.tex}
\usepackage{booktabs}
\usepackage{array}
\usepackage{url}

\title{Latent Fixture Reasoning}
\author{Anonymous Authors}
\date{}

\begin{document}
\maketitle

\begin{abstract}
Assembly robots often treat support relations and fixtures as visible scene
geometry or as constraints supplied before planning. This paper studies the
missing case: a part fails to move because an unobserved bolt, latch, adhesive
patch, or hidden stop is still constraining it. We propose \emph{latent fixture
reasoning}: infer a discrete hidden fixture cause from small manipulation
outcome residuals, then choose the release or reorientation action implied by
that cause. A bounded literature sweep over 1168 arXiv/Crossref
records suggests that the nearest work covers contact-rich manipulation,
constraint reasoning, tactile inference, active perception, and fixture-aware
assembly, but rarely treats hidden fixtures as an explicit action-selecting
state variable. In a deterministic toy benchmark, geometry-only manipulation
collides with many fixture modes, while latent fixture reasoning sharply
reduces wrong release actions. The evidence is modest, but it exposes a useful
state abstraction for robots that must learn what is holding the world in
place by acting on it.
\end{abstract}

\section{Introduction}
When a robot pulls on a part and it does not move, the failure is not just a
low-level control error. It can be evidence that something unseen is holding
the part in place. In assembly, repair, household manipulation, and field
robotics, that hidden cause may be a bolt, a latch, adhesive, a cable, a clamp,
or a geometric stop outside the camera view. A planner that represents only
visible geometry can repeatedly choose unsafe or unproductive actions.

We call this problem \emph{latent fixture reasoning}. The key object is an
interpretable hidden fixture variable. It is latent because the robot does not
observe the fixture directly. It is a fixture because it constrains motion. It
is reasoning, rather than only prediction, because the inferred cause changes
the next action: unbolt, unlatch, heat, reorient, or pull.

The paper makes three contributions. First, it separates latent fixture causes
from generic uncertainty in manipulation. Second, it defines a small posterior
interface from probe residuals to fixture hypotheses. Third, it provides a
minimal deterministic simulation showing why this state variable can reduce
wrong fixture actions.

\section{Problem}
Let $x$ denote the visible object state and let $z$ be an unobserved fixture
mode in $\{\mathrm{free}, \mathrm{bolt}, \mathrm{latch}, \mathrm{adhesive},
\mathrm{stop}\}$. A robot chooses a probe action $a_p$ and observes an outcome
residual $o$: displacement, lift compliance, twist compliance, and thermal
softening response. The robot then chooses a task action $a_t$ from direct pull,
unbolt, unlatch, heat-and-peel, or reorient. The central question is not whether
the robot can predict the next state in general. It is whether $o$ identifies a
fixture cause well enough to avoid the wrong physical release action.

\section{Mechanism}
The proposed controller maintains a small posterior over fixture modes,
\[
  p(z \mid o) \propto p(o \mid z)p(z).
\]
The likelihood model can be hand-written, learned, or calibrated from robot
data. In this paper we use a hand-written model only to isolate the mechanism.
The controller is allowed to perform a reversible probe before committing to a
release action. If confidence is low, it performs an additional reversible
probe rather than applying a high-force release action.

This differs from a generic world model in two ways. First, the latent variable
has a robotics meaning: it is a hidden support or fixture cause. Second, the
variable is directly coupled to an action menu. A high posterior on bolt selects
unbolt; a high posterior on adhesive selects heat-and-peel; a high posterior on
hidden stop selects reorientation.

\section{Evidence}
We implemented a deterministic toy benchmark with 2,000 seeds. Each seed
samples one fixture mode, generates noisy probe residuals, and evaluates three
policies. The geometry-only policy assumes visible geometry is complete. The
reactive policy first fails, then uses a single heuristic residual. The latent
fixture policy infers the fixture before choosing a release action.

\begin{table}[h]
\centering
\small
\begin{tabular}{lrrrrrr}
\toprule
Policy & Success & Coll. & Wrong & Steps & Return & Conf. \\
\midrule
Geometry only & 621 & 1379 & 1379 & 3.689 & -0.696 & 0.000 \\
Reactive & 1612 & 1767 & 388 & 5.758 & 0.332 & 0.207 \\
Latent fixture & 1995 & 5 & 5 & 7.077 & 0.871 & 0.941 \\
\bottomrule
\end{tabular}
\caption{Toy hidden-fixture benchmark. The mechanism is valuable when a
visible-geometry policy cannot distinguish free motion from a hidden fixture
constraint.}
\end{table}

The absolute numbers should not be read as benchmark performance. They show a
mechanism: when the hidden cause selects a different release action, estimating
that cause before committing can reduce physical mistakes.

\section{Related Work Boundary}
The literature sweep covers contact-rich manipulation, task-and-motion planning,
tactile inference, support-relation reasoning, robot world models, and assembly
fixture planning. These areas provide the hostile prior. The narrow boundary for
this paper is an interpretable latent fixture graph whose nodes are not merely
unobserved state, but action-selecting causes of constraint.

\section{Limitations}
The experiment is a toy model with hand-coded fixture types and synthetic
residuals. It does not prove that the abstraction scales to real assemblies.
The next required step is a physical or high-fidelity simulated benchmark where
fixtures are hidden from perception but revealed through controlled probes.

\section{Conclusion}
Latent fixture reasoning asks robots to treat failed motion as evidence about
what hidden support is still constraining the world. The contribution is a
small but useful shift in representation: from visible geometry plus generic
uncertainty to an interpretable hidden fixture cause that chooses the next safe
manipulation action.

\section*{References}
\small
The accompanying \texttt{docs/related\_work\_matrix.csv} contains the bounded recovery
literature sweep used for this mechanism note.

\end{document}
